\textbf{COLUNI 2025 -- Questão 02}\par\medskip Leia o texto abaixo e responda as questões de 1 a 3.\par\medskip

\begin{center}\textit{Meu Rosário}\end{center}

\noindent
\begin{minipage}[t]{0.49\linewidth}
Meu rosário é feito de contas negras e mágicas.\\
Nas contas de meu rosário eu canto Mamãe Oxum e falo\\
padres-nossos, ave-marias.\\
Do meu rosário eu ouço os longínquos batuques do\\
meu povo\\
e encontro na memória mal-adormecida\\
as rezas dos meses de maio de minha infância.\\
As coroações da Senhora, onde as meninas negras,\\
apesar do desejo de coroar a Rainha,\\
tinham de se contentar em ficar ao pé do altar\\
lançando flores.\\
As contas do meu rosário fizeram calos\\
nas minhas mãos,\\
pois são contas do trabalho na terra, nas fábricas,\\
nas casas, nas escolas, nas ruas, no mundo.\\
As contas do meu rosário são contas vivas.\\
(Alguém disse que um dia a vida é uma oração,\\
eu diria porém que há vidas-blasfemas).\\
Nas contas de meu rosário eu teço entumecidos\\
sonhos de esperanças.
\end{minipage}\hfill
\begin{minipage}[t]{0.49\linewidth}
Nas contas do meu rosário eu vejo rostos escondidos\\
por visíveis e invisíveis grades\\
e embalo a dor da luta perdida nas contas\\
do meu rosário.\\
Nas contas de meu rosário eu canto, eu grito, eu calo.\\
Do meu rosário eu sinto o borbulhar da fome\\
No estômago, no coração e nas cabeças vazias.\\
Quando debulho as contas de meu rosário,\\
eu falo de mim mesma em outro nome.\\
E sonho nas contas de meu rosário lugares, pessoas,\\
vidas que pouco a pouco descubro reais.\\
Vou e volto por entre as contas de meu rosário,\\
que são pedras marcando-me o corpo-caminho.\\
E neste andar de contas-pedras,\\
o meu rosário se transmuda em tinta,\\
me guia o dedo,\\
me insinua a poesia.\\
E depois de macerar conta por conto do meu rosário,\\
me acho aqui eu mesma\\
e descubro que ainda me chamo Maria.
\end{minipage}

\begin{center}(EVARISTO, Conceição. \textit{Poemas da recordação e outros movimentos}. Rio de Janeiro: Malê, 2017. p. 38−39)\end{center}

Leia o fragmento a seguir.

\par

Ao se aproximarem da porta de nossa casa, Crispina tombou no chão. Estava suja, tinha um cheiro ruim de suor, urina e flores mortas. Vi o horror se instaurar nos olhos de minha mãe. Não era a primeira, nem segunda, nem terceira vez que chegava alguém desvairado. E certamente não seria a última que se internaria em nossa casa, como diziam que faziam num hospital da capital para os que enlouqueciam. Não eram hóspedes, visitas ou convidados. Eram pessoas desconectadas de seu eu, desconhecidas de parentes e de si. Eram pessoas com encosto ruim, conhecidos e também desconhecidos de todos. Eram famílias que depositavam suas esperanças nos poderes de Zeca Chapéu Grande, curador de jarê, que vivia para restituir a saúde do corpo e do espírito aos que necessitavam. (VIEIRA JUNIOR, Itamar. \textit{Torto arado}. São Paulo: Todavia, 2019. p. 32−33)

\par

Considerando a relação entre o poema \textit{Meu Rosário} e o trecho da obra \textit{Torto arado}, no que se refere aos aspectos da constituição da religiosidade brasileira, marque a afirmativa INCORRETA.\begin{enumerate}
\item Em \textit{Meu Rosário} e em \textit{Torto arado}, ocorre vinculação entre memória, religiosidade e fazer poético, o que
       alude à coexistência da religiosidade europeia com a de origem africana na constituição da brasilidade.
\item O poema \textit{Meu Rosário} trata da religiosidade híbrida brasileira e da aceitação permanente do negro em
       uma cerimônia cristã: “As contas do meu rosário fizeram calos/ nas minhas mãos/ pois são contas do
       trabalho na terra, nas fábricas,/ nas casas, nas escolas, nas ruas, no mundo".
\item Em \textit{Torto arado}, Zeca Chapéu Grande é líder religioso da comunidade onde vive, atuando como um
       sacerdote do jarê, religião típica daquela região que mistura crenças africanas, católicas e indígenas.
\item O processo de religiosidade em \textit{Meu Rosário} é responsável pelo despertar da memória “mal-adormecida” do
       eu-poético e o faz recordar a infância e o preconceito: “as coroações da Senhora, onde as meninas negras,
       apesar do desejo de coroar a Rainha, tinham de se contentar em ficar ao pé do altar lançando flores”.\end{enumerate}
